\chapter{Переранжирование и реранкеры}\label{ch:reranking}

Информационный поиск в современных системах обработки естественного языка, как правило, организован в виде многоэтапного конвейера, в котором быстрый этап отбора кандидатов сменяется более точным, но вычислительно затратным этапом уточнения.
Этап переранжирования (reranking) играет ключевую роль в обеспечении высокого качества итоговой выдачи: он уточняет порядок документов, полученных на этапе первичного извлечения, используя более выразительные модели оценки релевантности.
В обзоре литературы (Раздел~\ref{sec:reranking}) систематизированы основные направления развития методов переранжирования; настоящая глава фокусируется на их практическом применении в двухэтапном конвейере и представляет авторский реранкер на основе компактной декодерной модели семейства GigaChat.

%% =================================================================
\section{Двухэтапный поиск и роль переранжирования}\label{sec:rr_two_stage}
%% =================================================================

Прямое вычисление релевантности запроса каждому документу коллекции с помощью выразительной модели невозможно при больших объёмах данных в силу вычислительных ограничений.
По этой причине применяется двухэтапная схема (Рисунок~\ref{fig:two_stage_retrieval}): на первом этапе быстрый извлекатель формирует множество кандидатов, на втором этапе реранкер переупорядочивает лишь это сокращённое множество.

\begin{figure}[htbp]
    \centering
    \resizebox{0.95\textwidth}{!}{
        \begin{tikzpicture}[
    font=\sffamily,
    node distance=1.3cm and 1.9cm,
    box/.style={
        draw,
        rounded corners,
        minimum width=2.8cm,
        minimum height=1.1cm,
        align=center,
        thick
    },
    db/.style={
        draw,
        cylinder,
        shape border rotate=90,
        aspect=0.3,
        minimum width=2.2cm,
        minimum height=1.6cm,
        align=center,
        thick
    },
    stage/.style={
        draw,
        dashed,
        rounded corners,
        inner sep=0.35cm
    },
    arr/.style={-{Stealth[length=3mm]}, thick}
]
    %% Запрос
    \node[box] (query) {Запрос\\пользователя};
    %% Стадия 1: отбор кандидатов
    \node[box, right=of query] (retriever) {Bi-encoder\\(GigaEmbeddings)};
    \node[db, above=of retriever] (index) {Индекс\\документов};
    %% Стадия 2: переранжирование
    \node[box, right=of retriever] (reranker) {Cross-encoder\\реранкер};
    %% Результат
    \node[box, right=of reranker] (result) {Top-$n$\\документов};

    %% Подписи стадий
    \node[stage, fit=(retriever)(index), label={[font=\sffamily\small]above:Отбор (recall)}] {};
    \node[stage, fit=(reranker), label={[font=\sffamily\small]above:Уточнение (precision)}] {};

    %% Стрелки
    \draw[arr] (query) -- (retriever);
    \draw[arr] (index) -- (retriever);
    \draw[arr] (retriever) -- node[above, font=\small] {top-$k$} (reranker);
    \draw[arr] (reranker) -- (result);
\end{tikzpicture}

    }
    \caption{Двухэтапная схема поиска: этап отбора кандидатов (bi-encoder) и этап уточнения (cross-encoder реранкер).}
    \label{fig:two_stage_retrieval}
\end{figure}

\textbf{Этап отбора (recall).} На первом этапе используется модель раздельного кодирования (bi-encoder), в которой запрос $q$ и документ $d$ кодируются независимо в векторы $\mathbf{h}_q$ и $\mathbf{h}_d$, а релевантность оценивается через косинусное сходство:
\begin{equation}
\label{eq:rr_biencoder}
    s_{\text{bi}}(q, d) = \cos(\mathbf{h}_q, \mathbf{h}_d) = \frac{\mathbf{h}_q^\top \mathbf{h}_d}{\|\mathbf{h}_q\|\,\|\mathbf{h}_d\|}.
\end{equation}
Поскольку представления документов вычисляются заранее (офлайн), поиск сводится к задаче приближённого поиска ближайших соседей (approximate nearest neighbor, ANN) в векторном пространстве, что обеспечивает обработку запросов за сублинейное время.
В качестве bi-encoder в настоящей работе используется модель GigaEmbeddings, описанная в Главе~\ref{ch:gigaembeddings}.

\textbf{Этап уточнения (precision).} Раздельное кодирование принципиально ограничено: запрос и документ не взаимодействуют до момента вычисления сходства, что препятствует моделированию тонких семантических связей между ними.
Реранкер устраняет это ограничение, выполняя совместное кодирование пары:
\begin{equation}
\label{eq:rr_crossencoder}
    s_{\text{cross}}(q, d) = f_\theta(q \oplus d),
\end{equation}
где $\oplus$ обозначает конкатенацию запроса и документа в единую последовательность, а $f_\theta$ --- модель, обеспечивающая полное взаимное внимание (full attention) над токенами запроса и документа.
Совместное кодирование существенно повышает точность оценки релевантности, однако требует отдельного прямого прохода модели для каждой пары $(q, d)$, что делает его применимым лишь к ограниченному числу кандидатов $k$ (как правило, $k \in [50, 200]$), отобранных на первом этапе.

%% =================================================================
\section{Архитектуры реранкеров}\label{sec:rr_architectures}
%% =================================================================

Можно выделить три основных класса архитектур переранжирования, различающихся способом моделирования взаимодействия между запросом и документом.

\textbf{Кросс-энкодеры.} Классический кросс-энкодер кодирует конкатенацию запроса и документа единой энкодерной моделью (например, на основе BERT~\cite{devlin2019bert}), после чего скалярная оценка релевантности извлекается из представления служебного токена.
Полное взаимное внимание обеспечивает наивысшую точность среди рассматриваемых классов ценой наибольших вычислительных затрат.

\textbf{Модели позднего взаимодействия.} Архитектуры типа ColBERT~\cite{2004.12832} занимают промежуточное положение: запрос и документ кодируются раздельно на уровне токенов, а сходство вычисляется через механизм мелкозернистого сопоставления (late interaction), что позволяет предвычислять представления документов при сохранении части преимуществ совместного кодирования (см. Раздел~\ref{subsec:late_interaction}).

\textbf{Реранкеры на основе больших языковых моделей.} Развитие генеративных LLM привело к появлению реранкеров, формулирующих оценку релевантности как задачу генерации.
Листовые (listwise) подходы, такие как RankZephyr~\cite{2312.02724}, упорядочивают список документов в рамках единого промпта.
Поточечные (pointwise) подходы оценивают каждую пару $(q, d)$ независимо, сводя задачу к бинарной классификации: модель оценивает вероятность того, что следующим сгенерированным токеном будет утвердительный ответ <<да>> (yes) против отрицательного <<нет>> (no).
Современная реализация данного подхода представлена в Qwen3~Embedding~\cite{2506.05176}, где реранкер строится поверх фундаментальной декодерной модели и демонстрирует передовые результаты.
Именно поточечная парадигма положена в основу авторского реранкера, описываемого далее.

%% =================================================================
\section{Реранкер на основе GigaChat-600M}\label{sec:rr_gigachat}
%% =================================================================

В настоящем разделе представлен авторский реранкер, построенный путём адаптации компактной декодерной модели GigaChat-600M семейства GigaChat~\cite{2506.09440} в поточечный кросс-энкодер.
В качестве основы используется базовая предобученная (pretrain) версия модели --- компактный плотный (dense) трансформер с приблизительно 600 миллионами параметров, без архитектуры смеси экспертов.
Выбор компактной модели обусловлен жёсткими требованиями к задержке на этапе уточнения, где реранкер применяется к десяткам или сотням кандидатов для каждого запроса.

\subsection{Формализация задачи}\label{subsec:rr_formulation}

Релевантность пары $(q, d)$ оценивается через вероятности утвердительного и отрицательного токенов, порождаемых моделью в ответ на инструктивный промпт.
Пусть $T(\cdot)$ --- шаблон, объединяющий инструкцию задачи $I$, запрос $q$ и документ $d$ в единую последовательность, а $z_{\text{yes}}$ и $z_{\text{no}}$ --- логиты соответствующих токенов на последней позиции.
Оценка релевантности определяется как нормированная вероятность утвердительного токена:
\begin{equation}
\label{eq:rr_score}
    s(q, d) = \frac{\exp(z_{\text{yes}})}{\exp(z_{\text{yes}}) + \exp(z_{\text{no}})}.
\end{equation}
Инструктивный шаблон $T(\cdot)$ задаётся следующим образом:
\begin{quote}
\ttfamily\small
Инструкция: \{$I$\}\\
Запрос: \{$q$\}\\
Документ: \{$d$\}\\
Релевантен ли документ запросу? Ответьте <<да>> или <<нет>>.
\end{quote}
Использование настраиваемой инструкции $I$ обеспечивает адаптируемость реранкера к различным определениям релевантности (например, поиск ответа, поиск по теме, поиск дубликатов) без изменения весов модели.

\subsection{Процедура обучения}\label{subsec:rr_training}

Обучение реранкера выполняется в режиме обучения с учителем (supervised fine-tuning, SFT).
Для каждого запроса $q$ формируются позитивный документ $d^+$ и множество жёстких негативных документов $\mathcal{N}_q$.
Жёсткие негативы добываются первой стадией поиска: для запроса извлекаются top-$k$ кандидатов с помощью GigaEmbeddings (Глава~\ref{ch:gigaembeddings}), из которых исключаются известные позитивы; такой подход обеспечивает обучение на трудноразличимых примерах, наиболее близких к распределению, наблюдаемому на этапе инференса.

Функция потерь представляет собой отрицательное логарифмическое правдоподобие корректной метки $l \in \{\text{yes}, \text{no}\}$:
\begin{equation}
\label{eq:rr_loss}
    \mathcal{L} = -\sum_{(q, d, l)} \log P\big(l \mid T(q, d)\big),
\end{equation}
где метка $l = \text{yes}$ соответствует позитивной паре $(q, d^+)$, а $l = \text{no}$ --- негативной паре $(q, n)$, $n \in \mathcal{N}_q$.
Следуя подходу Qwen3~Embedding~\cite{2506.05176}, этап масштабного слабоконтролируемого предобучения, применяемый при обучении моделей эмбеддингов, для реранкера не используется: совместное кодирование пары обеспечивает достаточно сильный обучающий сигнал уже на этапе дообучения с учителем.
Для повышения робастности к различным распределениям данных финальная модель может быть получена слиянием нескольких чекпойнтов дообучения посредством сферической линейной интерполяции (spherical linear interpolation, SLERP).

%% =================================================================
\section{Экспериментальная оценка}\label{sec:rr_evaluation}
%% =================================================================

\subsection{Протокол оценки}\label{subsec:rr_protocol}

Качество реранкера оценивается в рамках двухэтапной схемы: для каждого запроса извлекаются top-100 кандидатов с помощью bi-encoder GigaEmbeddings, после чего реранкер переупорядочивает полученное множество.
В качестве метрики используется нормированный дисконтированный накопленный выигрыш nDCG@10, стандартный для задач ранжированного поиска.
Оценка проводится на русскоязычных подмножествах ранжированного поиска бенчмарка ruMTEB~\cite{snegirev2025rumteb}, а также на классических наборах MS-MARCO~\cite{bajaj2018msmarco} и Natural Questions~\cite{kwiatkowski2019nq}.

\subsection{Базовые ориентиры}\label{subsec:rr_baselines}

В качестве ориентира для оценки качества используется семейство реранкеров Qwen3-Reranker~\cite{2506.05176}, демонстрирующее передовые результаты при сопоставимых масштабах моделей.
Результаты Qwen3-Reranker в протоколе переранжирования top-100 представлены в Таблице~\ref{tab:rr_qwen3}.
Модель Qwen3-Reranker-0.6B, сопоставимая по масштабу с авторским реранкером, превосходит реранкер BGE-reranker-v2-m3 того же размера на большинстве подмножеств (например, $+8.8$ пункта nDCG@10 на MTEB-R), что подтверждает эффективность поточечной парадигмы на основе декодерных LLM в компактном размерном классе.

\begin{table}[H]
    \centering
    \caption{Качество реранкеров Qwen3-Reranker (nDCG@10) в протоколе переранжирования top-100 кандидатов, отобранных bi-encoder Qwen3-Embedding-0.6B. Источник: \cite{2506.05176}.}
    \label{tab:rr_qwen3}
    \resizebox{\textwidth}{!}{
    \begin{tabular}{lccccccc}
        \toprule
        Модель & Парам. & MTEB-R & CMTEB-R & MMTEB-R & MLDR & MTEB-Code & FollowIR \\
        \midrule
        BGE-reranker-v2-m3              & 0.6B & 57.03 & 72.16 & 58.36 & 59.51 & 41.38 & $-0.01$ \\
        gte-multilingual-reranker-base  & 0.3B & 59.51 & 74.08 & 59.44 & 66.33 & 54.18 & $-1.64$ \\
        jina-multilingual-reranker-v2-base & 0.3B & 58.22 & 63.37 & 63.73 & 39.66 & 58.98 & $-0.68$ \\
        \midrule
        Qwen3-Reranker-0.6B     & 0.6B & 65.80 & 71.31 & 66.36 & 67.28 & 73.42 & 5.41 \\
        Qwen3-Reranker-4B       & 4B   & 69.76 & 75.94 & 72.74 & 69.97 & 81.20 & 14.84 \\
        Qwen3-Reranker-8B       & 8B   & 69.02 & 77.45 & 72.94 & 70.19 & 81.22 & 8.05 \\
        \bottomrule
    \end{tabular}
    }
\end{table}

\subsection{Результаты}\label{subsec:rr_results}

\todo[inline]{Вставить измеренные метрики реранкера GigaChat-600M (nDCG@10 на ruMTEB-R, MS-MARCO, NQ) и сравнение с базовыми ориентирами из Таблицы~\ref{tab:rr_qwen3}.}

Экспериментальная оценка авторского реранкера на основе GigaChat-600M проводится в соответствии с протоколом, описанным в Разделе~\ref{subsec:rr_protocol}.
Ожидается, что компактный реранкер обеспечит существенный прирост качества переранжирования относительно одноэтапного поиска с помощью bi-encoder, оставаясь при этом сопоставимым по задержке с моделями аналогичного масштаба.

%% =================================================================
\section{Выводы}\label{sec:rr_conclusions}
%% =================================================================

В настоящей главе рассмотрена роль переранжирования в двухэтапном конвейере информационного поиска и представлен авторский реранкер на основе компактной декодерной модели GigaChat-600M.
Реранкер построен по поточечной парадигме: оценка релевантности извлекается из вероятностей утвердительного и отрицательного токенов в ответ на инструктивный промпт, а обучение выполняется в режиме дообучения с учителем на жёстких негативах, добываемых моделью GigaEmbeddings.
Совместно с моделью эмбеддингов (Глава~\ref{ch:gigaembeddings}) реранкер образует основу двухэтапного поискового стека, используемого в системе генерации с извлечением информации (Глава~\ref{ch:rag}) и в конвейере обеспечения безопасности (Глава~\ref{ch:safety}).
